\documentclass[a4paper,9pt]{article}

\usepackage{fullpage}
\usepackage[all]{xy}
\usepackage[utf8]{inputenc}
\usepackage{amsmath,amsthm}
\usepackage{amsfonts}
%\usepackage{algorithm}
\usepackage{graphicx}
%\usepackage[compatible,noend]{algpseudocode}
\usepackage{latexsym}
\usepackage{float,color}
\usepackage{xcolor}
\usepackage{framed}


\definecolor{shadecolor}{RGB}{0,0,0}


\newenvironment{terminal}{\color{white}\begin{snugshade*}}
{\end{snugshade*}\color{black}}


%%%%%%%%%%%%%%% COMMANDES %%%%%%%%%%%%%%%%%%%%%%%%%%%

\newcommand{\itemb}{\item[$\bullet$]}

%% Les ensembles de modelisation
\newcommand{\Ki}{\mathcal{K}}
\newcommand{\Ci}{\mathcal{C}}
\newcommand{\Mi}{\mathcal{M}}
\newcommand{\Ei}{\mathcal{E}}
\newcommand{\Di}{\mathcal{D}}

\newcommand{\Fd}{\mathbb{F}}
\newcommand{\Znat}{\mathbb{Z}}
\newcommand{\Nnat}{\mathbb{N}}
\newcommand{\pgcd}{\mbox{pgcd}}
%%%%%%%%%%%%%%% ENVIRONEMENTS %%%%%%%%%%%%%%%%%%%%%%%


\theoremstyle{plain}
\newtheorem{lemme}{Lemme}
\newtheorem{algorithme}{ Algorithme }
\newtheorem{corolaire}{Corollaire}
\newtheorem{propriete}{Propri\'et\'e}
\newtheorem{proposition}{Proposition}
\newtheorem{theoreme} {Th\'eor\`eme}
\newtheorem{definition}{D\'efinition}

\theoremstyle{definition}
\newtheorem*{probleme}{Problème}
\newtheorem*{cout}{Coût}
\newtheorem*{fait}{Fait}
\newtheorem*{notation}{Notation}
\newtheorem*{complexite}{Complexit\'e}
\newtheorem{exercice}{Exercice}
%\theoremstyle{remark}
\newtheorem{remarque}{Remarque}
\newtheorem{exemple}{Exemple}


%%%%%%%%%%%%%% DOCUMENTS %%%%%%%%%%%%%%%%%%%%%%%%%%%%


\begin{document}


\hbox{\raisebox{0.4em}{\vrule depth 0pt height 0.5pt width \textwidth}}
\noindent \textbf{Université de Perpignan} \hfill  Année \the\year \\
Licence 3 Informatique \hfill C. Legrand \\

\begin{center}
  {  {\scshape \large Contrôle Terminal 1} \\
    \vrule depth 0pt height 0.5pt width 1cm
    \medskip
    
{\scshape    S\'ecurit\'e Informatique}  \\
}
\end{center}
\hbox{\raisebox{0.4em}{\vrule depth 0pt height 0.5pt width \textwidth}}

\bigskip

\bigskip


\noindent\textit{Recommandations :}
\begin{itemize}
\item Durée: 3h.
\item Notes de cours et de TD autorisées.
  \item Documents du moodle autorisés. 
  \end{itemize}
\noindent\textit{Commandes bash utiles :} \verb|head|, \verb|tail| et \verb|cut|. Pour \verb|cut| on définit le délimiteur avec l'option \verb|-d| et le numéro du champs \verb|n| sélectionné avec \verb|-fn|.  Voici un exemple où l'on récupère le second champs:
\begin{small}
\begin{terminal}
\begin{verbatim}
clegrand@mon-pc:~$  cat fichier.txt 
nom prenom age
Durand Michel 19
Esposito Anna 35
Dufour Elsa 47
clegrand@mon-pc:~$ cat fichier.txt | cut -d " " -f2
prenom
Michel
Anna
Elsa
\end{verbatim}
\end{terminal}
\end{small}
%\begin{itemize}
%\item Tout document (papier ou électronique (pdf, doc, txt) ) est  autorisé.
%\item Vous devrez rendre une copie \'electronique (en format .txt) de l'examen.
%\item Vous l'enverrez en pi\`ece jointe à l'adresse email suivante
%\begin{center}
%\texttt{cedric.legrand@univ-perp.fr}
%\end{center}
%ou me la transmettrez sur une clef usb.
%\end{itemize}

\begin{exercice}[Parcours de fichier ligne à ligne]
  \begin{enumerate}
\item Faire un script \verb|bash| qui prend en argument un nom de fichier \verb|F| et un entier \verb|n|. Le script affiche sur la sortie standard la \verb|n|-ième ligne du fichier \verb|F| (utilisez \verb|head| et \verb|tail|). %  la commande \verb|bash| qui à partir d'un fichier \verb|f| affich
\item Faire un script qui prend un fichier en argument et affiche sur la sortie standard chaque ligne précédée de \verb|ligne <numero-ligne>:|. Voici une illustration (le scrit de cette question doit agir comme \verb|exo1-q2.sh|):
  \begin{small}
\begin{terminal}
\begin{verbatim}
clegrand@mon-pc:~$ cat chiffre.txt
zero
un
deux
trois
quatre
cinq
six
sept
huit
neuf
clegrand@mon-pc:~$ ./exo1-q2.sh chiffre.txt
ligne 1: zero
ligne 2: un
ligne 3: deux
ligne 4: trois
ligne 5: quatre
ligne 6: cinq
ligne 7: six
ligne 8: sept
ligne 9: huit
ligne 10: neuf
clegrand@clegrand-OptiPlex-990:~/Desktop/Securite/Exam2022/CT1$ 
\end{verbatim}
\end{terminal}
\end{small}

  \end{enumerate}
\end{exercice}

%\bigskip
\newpage

\begin{exercice}[Virus bash par entrelacement]
%// Virus script bash par entrelacement (faire un mini exemple)

Le but de cet exercice est d'écrire le  code en script bash d'un virus qui va disséminer son code dans le code hote de la manière suivante:
\begin{center}
\includegraphics[width=7cm]{entrelacement} 
\end{center}

\begin{enumerate}
\item Ecrire le  un script bash \texttt{entrelacement.sh} qui prend deux fichiers (\texttt{cible.sh} et \texttt{codeviral.sh}) en argument et qui fabrique un troisième fichier \texttt{resultat.sh} comme suit:
\begin{itemize}
\item si le nombre de ligne de \texttt{cible.sh} est inférieur à \texttt{codeviral.sh} alors ne rien faire.
\item si le nombre de ligne de \texttt{cible.sh} est supérieur à \texttt{codeviral.sh} alors définir 
  $$s=\left\lfloor\frac{\mbox{nombre de ligne de \texttt{cible.sh}}}{\mbox{nombre de ligne de \texttt{codeviral.sh}}}\right\rfloor.$$
  Rappel : $\lfloor x \rfloor = $ partie entière de $x$.
\item remplir le fichier \texttt{resultat.sh} en ajoutant alternativement $s$ lignes de  \texttt{cible.sh} suivies d'une ligne de \texttt{codeviral.sh}.
\end{itemize}

%\noindent\textit{Indication.} Vous pourrez utiliser la boucle suivante
%\begin{verbatim}
%while read ligne
%do
%
%done < $file
%\end{verbatim}
%qui lit le fichier \texttt{\$file} ligne par ligne.

\item Faire un script  \texttt{restauration.sh} qui prend en argument  le nombre \verb|s| et  un fichier \texttt{resultat.sh} fabriqué dans la question précédente. Le script  restaure les  deux fichiers (\texttt{cible.sh} et \texttt{codeviral.sh}). 
\medskip



\item Combiner les deux scripts \texttt{entrelacement.sh} et \texttt{restauration.sh} pour fabriquer votre virus qui infectera uniquement le répertoire courant:
  \begin{itemize}
  \item Le code de restauration restaure le code hote et le code viral dans deux fichiers temporaires, il exécute d'abord le code viral, puis il exécute le code hote.
  \item Le code viral parcourt les fichiers \verb|*.sh| du répertoire courant et ajoute son code par entrelacement comme dans le question 1).
    \end{itemize}
\item Expliquez une stratégie possible pour rendre la segmentation dans la partie \texttt{entrelacement.sh} aléatoire : les nombre de lignes du code hôte consécutive ne sont plus fixes (égaux à \verb|s|) mais varient de façon aléatoire.
\end{enumerate}
\end{exercice}


\end{document}




\begin{exercice}
Tentative d'attaque du ver \texttt{Code Red}: les logs d'un serveur web contient une série d'entrées du type [\texttt{adresse source, date, requête, résultat, octet transférer}]. Expliquez la particularité du log suivant:
\begin{verbatim}
64.163.156.186-[10/Feb/2005:02:37:24 +0100] << GET /default.ida ? NNNNNNNNNNNNNNNNNN
NNNNNNNNNNNNNNNNNNNNNNNNNNNNNNNNNNNNNNNNNNNNNNNNNNNNNNNNNNNNNNNNNNNNNNNNNNNNNNNNNNNN
NNNNNNNNNNNNNNNNNNNNNNNNNNNNNNNNNNNNNNNNNNNNNNNNNNNNNNNNNNNNNNNNNNNNNNNNNNNNNNNNNNNN
NNNNNNNNNNNNNNNNNNNNNNNNNNNNNNNNNNNNNNNNNNNNNNNNNNNNNNNNNNNNNNNNNNNNNNNNNNNNNNNNNNNN
NNNNNNNNNNNNNNNNNNNNNNNNNNNNNNNNNNNNNNNNN%u9090%u6858%ucbd3%u7801%u9090%u6858%ucbd3%
u7801%u9090%u6858%ucb%u78011%u90901%u9090%8190%u00c3%u0003%u8b00%531b%u53ff%u0078%u0
000%u00=a HTTP/1.0" 400 333
\end{verbatim}

\end{exercice}


%% la suite N sert a priori de bourrage 
%% pour que la fin de la chaine deborde du tampon
%% et qu'un code malicieux soit executé à la place 
%% du code du serveur web

\begin{exercice}
Soit le site \texttt{www.unsitecomplice.fr}, le moyen de vérifier qu'il est vulnérable à une attaque \textit{cross site scripting} est de demander l'affichage d'une page inexistante. La particularité du nom de cette page est le fait qu'il contient des balises HTML:
\begin{verbatim}
   http://wwww.unsitecomplice.fr/<B>nimportequoi</B>&g.html
\end{verbatim}
Erreur la page nimportequoi.html est introuvable.

\begin{enumerate}
\item Pour quelles raisons le site est-il vulnérable? 
\item Qui peut-être la victime : les internautes ou le site web?
\item Donner un scénario d'attaque possible via un \texttt{www.unsitecomplice.fr} (par exemple en utilisant l'ingénierie sociale).
\end{enumerate}
\end{exercice}

%% il est vulnerable car les balises html du nom de fichier ont ete interprete
%% le site web peut servir de ``complice''. 
%% Par ingenieurie sociale on redirige vers le site en metant dans la barre 
%% d'adresse un faux fichier contenant une fausse page html
%% et un script pour pouvoir recuperer, par exemple, un mot de passe 


%// Lecture d'un virus du giant black book

%// Virus/programmation systeme (l'idée etait quoi deja???)

\begin{exercice}
Modifier le virus \texttt{vcomp\_ex}  vu en cours afin que celui-ci exécute le code viral et le code cible en parallèle. 
\medskip

\noindent\textit{Indication.} Vous pourrez utiliser la fonction 
\begin{verbatim}
int pthread_create(pthread_t * thread, 
                   pthread_attr_t * attributs,
                   void * (* fonction)(void * argument),
                   void *argument); 
\end{verbatim}
de \texttt{<pthread.h>}.
Discuter ensuite des avantages et inconvénients de cette approche.
\end{exercice}

\begin{exercice}
Modifier le virus résident vu en TD (exercice 1, Feuille TD 4 sur les vers), pour qu'il change de nom régulièrement:
\begin{enumerate}
\item Première approche: vous modifierez le nom en ajoutant  le \texttt{pid} (du père ou du fils) dans le nom du programme.
\item Deuxième approche: vous analyserez les processus en cours d'exécution et utiliserez de façon aléatoire le nom de l'un d'eux  pour le nom du programme lors du rafraîchissement.
\end{enumerate}
\end{exercice}

